\section*{Bloque A. Fundamentos teóricos y conceptuales}

\noindent\textit{Preguntas 1 a 30. Evalúan la comprensión del equilibrio general, la eficiencia de Pareto, la caja de Edgeworth, la curva de contrato, el núcleo y los teoremas fundamentales del bienestar.}

\q{El equilibrio general se define como:}
\begin{opts}
  \item el equilibrio de un solo mercado, sin interacción con los demás.
  \item una situación en la que solo los consumidores maximizan su utilidad.
  \item una asignación fijada por el regulador.
  \item una situación en la que todos los mercados se vacían simultáneamente y los precios se ajustan.
\end{opts}

\q{Un equilibrio walrasiano es:}
\begin{opts}
  \item cualquier asignación Pareto eficiente.
  \item un vector de precios tal que la oferta iguala la demanda en todos los mercados.
  \item un equilibrio que solo considera el lado de la producción.
  \item un precio fijado por el Estado.
\end{opts}

\q{En equilibrio general, solo importan los precios relativos porque:}
\begin{opts}
  \item el dinero es neutral en el corto plazo.
  \item los precios absolutos siempre son cero.
  \item las demandas y ofertas dependen de los cocientes de precios y se puede elegir un numerario.
  \item el Estado fija todos los precios.
\end{opts}

\q{La ley de Walras establece que:}
\begin{opts}
  \item la suma del valor de los excesos de demanda es cero a los precios vigentes.
  \item todos los mercados tienen exceso de demanda positivo.
  \item los precios son iguales en todos los mercados.
  \item la oferta siempre es igual a la demanda en cada mercado.
\end{opts}

\q{La función de exceso de demanda es homogénea de grado:}
\begin{opts}
  \item 1 en precios.
  \item 2 en precios.
  \item $-1$ en precios.
  \item 0 en precios.
\end{opts}

\q{La existencia de un equilibrio competitivo requiere, entre otros supuestos:}
\begin{opts}
  \item externalidades negativas.
  \item preferencias convexas y funciones de demanda continuas.
  \item poder de mercado de una empresa.
  \item precios rígidos fijados por el gobierno.
\end{opts}

\q{La principal diferencia entre el equilibrio parcial y el equilibrio general es que:}
\begin{opts}
  \item el parcial ignora el mercado del bien analizado.
  \item el general estudia un solo mercado.
  \item el general incorpora la interdependencia entre todos los mercados.
  \item el parcial supone competencia imperfecta.
\end{opts}

\q{Una asignación es Pareto eficiente si:}
\begin{opts}
  \item todos los consumidores reciben la misma cantidad.
  \item se maximiza la producción total.
  \item no existe comercio entre los agentes.
  \item no es posible mejorar a alguien sin empeorar a otro.
\end{opts}

\q{Una mejora de Pareto es:}
\begin{opts}
  \item un cambio en el que alguien mejora y otro empeora.
  \item un cambio en el que todos empeoran.
  \item una decisión adoptada por el regulador.
  \item un cambio en el que al menos uno mejora y nadie empeora.
\end{opts}

\q{El Primer Teorema del Bienestar establece que:}
\begin{opts}
  \item todo equilibrio competitivo es Pareto eficiente bajo ciertos supuestos.
  \item toda asignación eficiente es equitativa.
  \item los mercados siempre fallan.
  \item el gobierno debe fijar los precios.
\end{opts}

\q{El Primer Teorema del Bienestar puede fallar cuando:}
\begin{opts}
  \item hay competencia perfecta.
  \item existen externalidades, bienes públicos, poder de mercado o información asimétrica.
  \item los mercados son completos.
  \item todos los agentes son precio-aceptantes.
\end{opts}

\q{El Segundo Teorema del Bienestar establece que:}
\begin{opts}
  \item todo equilibrio competitivo es equitativo.
  \item no se puede alcanzar la eficiencia sin monopolio.
  \item toda asignación Pareto eficiente puede sostenerse como equilibrio competitivo con transferencias lump-sum, bajo ciertos supuestos.
  \item la eficiencia exige precios fijados por el gobierno.
\end{opts}

\q{El Segundo Teorema del Bienestar utiliza transferencias lump-sum porque:}
\begin{opts}
  \item no distorsionan las decisiones marginales y solo redistribuyen ingreso.
  \item siempre aumentan la eficiencia.
  \item las pagan las empresas.
  \item requieren precios máximos.
\end{opts}

\q{La caja de Edgeworth representa:}
\begin{opts}
  \item la frontera de posibilidades de producción.
  \item las asignaciones de dos bienes entre dos consumidores para dotaciones totales dadas.
  \item solo los precios de equilibrio.
  \item el presupuesto del regulador.
\end{opts}

\q{El punto de dotación en la caja de Edgeworth corresponde a:}
\begin{opts}
  \item la curva de contrato.
  \item la asignación inicial de recursos.
  \item el equilibrio de mercado.
  \item el óptimo social.
\end{opts}

\q{La curva de contrato es:}
\begin{opts}
  \item el conjunto de asignaciones Pareto eficientes.
  \item la recta presupuestaria.
  \item la curva de demanda del mercado.
  \item la función de producción.
\end{opts}

\q{El núcleo o core de una economía es:}
\begin{opts}
  \item el conjunto de asignaciones que ninguna coalición puede mejorar.
  \item el conjunto de todas las asignaciones factibles.
  \item el conjunto de las asignaciones igualitarias.
  \item el conjunto de asignaciones fijadas por el gobierno.
\end{opts}

\q{En la caja de Edgeworth, el área en forma de lente representa:}
\begin{opts}
  \item asignaciones ineficientes.
  \item únicamente la dotación inicial.
  \item la curva de contrato.
  \item asignaciones que mejoran a ambos consumidores respecto de la dotación.
\end{opts}

\q{En un óptimo de Pareto interior de intercambio puro se cumple:}
\begin{opts}
  \item $MRS_A>MRS_B$.
  \item $MRS_A<MRS_B$.
  \item $MRS=0$.
  \item $MRS_A=MRS_B$.
\end{opts}

\q{En una economía con producción, la eficiencia requiere que:}
\begin{opts}
  \item $MRS=0$.
  \item $MRS=MRT$.
  \item $MRS>MRT$.
  \item $P=CMe$.
\end{opts}

\q{La frontera de posibilidades de utilidad muestra:}
\begin{opts}
  \item los costos de producción.
  \item los precios de equilibrio.
  \item la máxima utilidad de un consumidor dada la utilidad del otro.
  \item la dotación inicial de recursos.
\end{opts}

\q{Una función de bienestar social:}
\begin{opts}
  \item describe solo las preferencias individuales.
  \item agrega las utilidades individuales en un ordenamiento social.
  \item siempre es Pareto eficiente.
  \item ignora por completo la equidad.
\end{opts}

\q{Una asignación Pareto eficiente:}
\begin{opts}
  \item no necesariamente es equitativa.
  \item siempre es igualitaria.
  \item maximiza la igualdad.
  \item exige que las utilidades sean iguales.
\end{opts}

\q{Una transferencia lump-sum:}
\begin{opts}
  \item distorsiona los precios relativos.
  \item es un impuesto al consumo.
  \item redistribuye ingreso sin crear distorsiones marginales.
  \item es un precio máximo.
\end{opts}

\q{El criterio de Kaldor-Hicks:}
\begin{opts}
  \item exige que nadie empeore.
  \item es idéntico a una mejora de Pareto.
  \item ignora la eficiencia.
  \item considera una mejora potencial de Pareto si los ganadores pueden compensar a los perdedores.
\end{opts}

\q{Una externalidad negativa genera:}
\begin{opts}
  \item ineficiencia asignativa y una falla del Primer Teorema del Bienestar.
  \item eficiencia automática.
  \item equidad.
  \item precios competitivos.
\end{opts}

\q{Un bien público puro se caracteriza por:}
\begin{opts}
  \item no rivalidad y no exclusión.
  \item rivalidad y exclusión.
  \item ser producido por un monopolio.
  \item tener un precio positivo.
\end{opts}

\q{Un monopolio genera pérdida de bienestar porque:}
\begin{opts}
  \item fija $P=CMg$.
  \item produce más que la competencia.
  \item fija $P>CMg$ y produce menos que la cantidad competitiva.
  \item no tiene costos.
\end{opts}

\q{El Primer Teorema del Bienestar requiere, entre otros supuestos:}
\begin{opts}
  \item mercados completos, competencia perfecta, información perfecta y ausencia de externalidades.
  \item monopolio.
  \item externalidades.
  \item precios rígidos.
\end{opts}

\q{Si los mercados son incompletos o hay información asimétrica:}
\begin{opts}
  \item el Primer Teorema se cumple siempre.
  \item la eficiencia aumenta.
  \item el equilibrio competitivo puede no ser Pareto eficiente.
  \item no hay intercambio.
\end{opts}
